\documentclass[10pt,a4paper]{amsart}
\usepackage[margin=19mm]{geometry}
\usepackage{amsmath,amssymb,mathtools}
\usepackage[scr=rsfs,cal=euler]{mathalfa}
\usepackage{xfrac}
\usepackage[unicode,hidelinks,bookmarks=false]{hyperref}
\newcommand{\ie}{i.e.\ }
\newcommand{\cf}{cf.\ }
\newcommand{\resp}{respectively\ }
\newcommand{\Subsection}{\subsection}
\providecommand{\nobreakdash}{\nobreak\hbox{-}}
\setlength{\emergencystretch}{2em}
\multlinegap0pt
\usepackage{bm}
\usepackage{bbm}
\usepackage{dsfont}
\usepackage{xspace}
\usepackage{cancel}
\usepackage{mathtools}
\usepackage{cleveref}
\usepackage[shortlabels]{enumitem}
\usepackage{amsthm}
\makeatletter
\@ifundefined{definition}{\newtheorem{definition}{Definition}}{}
\@ifundefined{proposition}{\newtheorem{proposition}{Proposition}}{}
\makeatother
\makeatletter
\newcommand{\R}{\mathbb{R}}
\newcommand{\aep}{a_\e}
\newcommand{\be}{b_\e}
\newcommand{\e}{\epsilon}
\newcommand{\fpe}{f\star\phi_\e}
\expandafter\ifx\csname fromtex\endcsname\relax
\newenvironment{fromtex}{}
% {\par\color{blue}}{\par}
\fi
\newcommand{\gep}{{\gamma_\epsilon}}
\makeatother
\title[Unbalanced Minimal Matchings for Measurable Scale Invariant ]{Unbalanced Minimal Matchings for Measurable Scale Invariant Costs on $\mathbb{R}^d$}
\author{Omer Angel, Nitya Gadhiwala}
\date{Source: arXiv:2609.19426v1 (CC BY 4.0)}
\begin{document}
\maketitle

\noindent\emph{Source and licence.} Unbalanced Minimal Matchings for Measurable Scale Invariant Costs on $\mathbb{R}^d$. Version arXiv:2609.19426v1 (CC BY 4.0), distributed under the Creative Commons Attribution 4.0 International licence, as recorded in the arXiv licence metadata for this exact version and shown on the abstract page https://arxiv.org/abs/2609.19426v1. Full rights notice: licenses/arxiv-2609.19426-CC-BY-4.0.txt.\par\smallskip

\noindent\textit{Editorial scope.} Counted unit: the Proposition whose source label is \texttt{recoverf} in the article (source0.tex lines 1594--1605, source label \texttt{recoverf}), delivered together with the article's own complete original proof of it (source0.tex lines 1607--1698, one \texttt{proof} environment of 92 lines). No mathematics is written, repaired or re-derived: statement and proof are the article's own text line for line. Required context retained uncounted: thm:monotone (lines 1016--1019); thm:continuous (lines 1265--1269); scinvfrommin (lines 1517--1529); approxid (lines 1541--1549). Every definition, display and label of those lines is retained unchanged. Omitted from this package and not redistributed: 1--1015, 1020--1264, 1270--1516, 1530--1540, 1550--1593, 1606--1606, 1699--3353. Nothing inside a retained excerpt is omitted or altered: every retained body is delivered exactly as the article's lines are written, so no omission receipt of any kind accompanies this package. Citations: the retained excerpts carry no citation command at all, so no bibliography entry of the article is transcribed and no reference is redistributed. Reference adaptation: none was needed and none was made; every reference inside the delivered unit resolves inside it, so no unresolved reference or citation is produced. Printed slips kept literal: no printed word, symbol, hypothesis or number of the counted statement or of its proof was altered. Layout and class: the article's own class is \texttt{article} with packages including amsmath, amssymb, amsfonts, amstext, amsthm, url, xspace, color, enumerate, verbatim, subcaption, graphicx, hyperref, cleveref; the wrapper is the lane's pinned \texttt{amsart} editorial wrapper at 10pt on A4, which restates the theorem-like environments the retained text needs and adds the article's own macro definitions that the retained text needs (\texttt{\textbackslash R}, \texttt{\textbackslash aep}, \texttt{\textbackslash be}, \texttt{\textbackslash e}, \texttt{\textbackslash fpe}, \texttt{\textbackslash gep}), with extra wrapper packages (bm, bbm, dsfont, xspace, cancel, mathtools, cleveref, enumitem). The delivered numbering is the unit's own. Assets: no retained span contains a figure environment, a graphics-inclusion command, a PDF-page inclusion, a table or a TikZ picture; no image file is copied into this package, the package declares no asset dependency and no image file is redistributed with it. Licence: version arXiv:2609.19426v1 of this article is distributed under the Creative Commons Attribution 4.0 International licence, as recorded in the arXiv licence metadata for this exact version and shown on the abstract page https://arxiv.org/abs/2609.19426v1. A full rights notice is delivered at licenses/arxiv-2609.19426-CC-BY-4.0.txt. Characters: every retained excerpt is pure ASCII, so the delivered engine, XeLaTeX with its default Latin Modern text font, reports no missing character.
\begin{proposition}[monotone]\label{thm:monotone}
  A measurable scale invariant cost function defined along a ray $f:(0,\infty)\rightarrow \mathbb{R}$ must be
  monotone.
\end{proposition}

\begin{proposition}\label{thm:continuous}
  A monotone, measurable scale invariant cost function on
  $\mathbb{R}^d$ can't have a discontinuity of size $k$ for any $k>
  0$.
\end{proposition}

\begin{proposition}\label{scinvfrommin}
  Let $g:(0,\infty)\rightarrow\R$ be a scale invariant cost function that is smooth,
  non-constant and non-decreasing. Then,
   there exist
$a,b,\gamma\in\mathbb{R}$, such that
\begin{equation}\label{eq:cost1dim}
  g(x) =
  \begin{cases}
    ax^\gamma +b, & \gamma\neq 0, \qquad \mbox{or} \\
    a\log x + b, & \gamma=0.
  \end{cases}
\end{equation}
\end{proposition}

\begin{definition}[approximate identity]\label{approxid}
  An approximate identity is a family of functions
  $\{\phi_\epsilon\}_{0<\epsilon<1}\in C^\infty_c(0,\infty)$ such that
  \begin{enumerate}[(i)]
  \item $\int_0^\infty \phi_\epsilon(a) da =1$ for all $\epsilon>0$,
  \item $\phi_\epsilon(x)\geq 0$ for all $x,\epsilon>0$,
    \item $\mbox{supp}(\phi_\epsilon) = (1-\epsilon,1+\epsilon)$.
  \end{enumerate}
\end{definition}

\begin{proposition}[Recovering $f$ from the convolution]\label{recoverf}
  Assume $f:(0,\infty)\rightarrow\mathbb{R}$ is a measurable scale
  invariant cost function along a ray. Then, there exist
  $a,b,\gamma\in\mathbb{R}$ such that
  \[
f(x) =
\begin{cases}
  ax^\gamma +b,& \gamma\neq 0, \qquad \mbox{or} \\
  a\log x +b,& \gamma =0.
\end{cases}
\]
\end{proposition}

\begin{proof}
  We established earlier that $f$ must be monotone
  (\cref{thm:monotone}), continuous (\cref{thm:continuous}), and
  the pointwise limit of the convolution above
  (\cref{approxid}).
  We know from \cref{scinvfrommin} that for each $\epsilon>0$, there exist
  $a_\e,b_\e,\gamma_\e\in\mathbb{R}$ such that
  \[
\fpe(x) =
\begin{cases}
  a_\e
  x^{\gamma_\e} +b_\e,& \gamma_e\neq 0, \qquad \mbox{or} \\
  a_\e\log x +b_\e, & \gamma_\e =0,
\end{cases}
\]
and $f\star\phi_\epsilon(x)\rightarrow f$ pointwise. We consider the
following two cases.


{\bf Case 1:} There are infinitely many $\e>0$ , such that
$\fpe(x) =a_\e \log x + b_\e$.
For simplicity, we index this sequence by $\e$. Due to pointwise
convergence of this sequence $\fpe$ to $f$ for all $x>0$, we have
\begin{align}
  f(1)& = \lim_{\e\rightarrow 0} \fpe(1) =  \lim_{\e\rightarrow 0} a_\e
  \log 1 +b_\e =  \lim_{\e\rightarrow 0} b_\epsilon, \label{eq:f1} \\
  f(e)& = \lim_{\e\rightarrow 0} \fpe(e) =  \lim_{\e\rightarrow 0} a_\e
  \log e +b_\e =  \lim_{\e\rightarrow 0}a_\e + b_\epsilon. \label{eq:fe}
\end{align}
 Now, \cref{eq:f1} implies $ \lim_{\e\rightarrow 0} b_\epsilon $ exists and equals
 $f(1)$. This, in addition to \cref{eq:fe} implies
 $\lim_{\e\rightarrow 0} a_\e $ also exists and equals $f(e)-f(1)$.
 Since the subsequence $\fpe(x) = \aep \log x +\be$ converges to
 $f(x)$ pointwise and both $\aep$ and $\be$ converge, we conclude that
 for $a = f(e)-f(1)$ and $b = f(1)$, we have
 \[f(x) = a\log x +b. \]

 {\bf Case 2:} There are infinitely many $\epsilon>0$ such that
 $\fpe=a_\e x^{\gamma_e} + b_\e$. A similar argument to the one above
 can be used to show that $\lim_{\e\rightarrow 0}a_\e$,
 $\lim_{\e\rightarrow 0}b_\e$ and $\lim_{\e\rightarrow 0}\gamma_\e$ all exist and that $f$ must equal
 $a x^\gamma +b$ where $a$, $b$ and $\gamma$ are the respective limits.
 \iffalse
This subsequence (which we index using $\e$ for simplicity) also converges to
$f$ pointwise. So, we have
\begin{align*}
  f(1) & = \lim_{\e\rightarrow 0} \fpe(1) = \lim_{\e\rightarrow 0}
  \aep+\be \\
  f(2) &  =  \lim_{\e\rightarrow 0} \fpe(2) = \lim_{\e\rightarrow 0}
         \aep 2^\gep +\be \\
  f(4) & =  \lim_{\e\rightarrow 0} \fpe(4) = \lim_{\e\rightarrow 0}
  \aep 4^\gep +\be
\end{align*}
Since the above limits all exist, we can get
\begin{align} \label{eq:f21}
  f(2)-f(1) = & \lim_{\e\rightarrow 0} \aep (2^\gep-1) \\
  f(4)-f(2) = &
\lim_{\e\rightarrow 0} \aep  2^\gep(2^\gep -1)
\end{align}
Since the above two limits exist and have non-zero limits, we can
divide them (taking a subsequence if necessary) to conclude that
$\lim_{\e\rightarrow 0} \gep$ exists.
\[
\lim_{\e\rightarrow 0} 2^\gep = \frac{f(4)-f(2)}{f(2)-f(1)}
\implies  \lim_{\e\rightarrow 0} \gep = \log_2\left(
  \frac{f(4)-f(2)}{f(2)-f(1)} \right)
\]
Now, since $\lim_{\e\rightarrow 0} 2^\gep$ exists, we know
$\lim_{\e\rightarrow 0} 2^\gep -1$ also exists. Combining this with
\cref{eq:f21}, we get
\begin{align*}
  %\label{eq:a}
  f(2)-f(1) = & \lim_{\e\rightarrow 0}\aep \left(
                \frac{f(4)-f(2)}{f(2)-f(1)} \right)\\
  \implies \qquad  \lim_{\e\rightarrow 0}\aep = &
  \frac{(f(2)-f(1))^2}{f(4)+f(1)-2f(2)}.
\end{align*}
Now that we know $\lim_{\e\rightarrow 0} \aep$ exists, we can say that
$\lim_{\e\rightarrow 0} \be = f(1) - \lim_{\e\rightarrow 0} \aep$.
Since limits exist for $\aep,\be,\gep$ as $\e$ approaches 0, we know
that for
\begin{align*}
  a &= \lim_{\e\rightarrow 0} \aep =
  \frac{(f(2)-f(1))^2}{f(4)+f(1)-2f(2)} \\
  b & = \lim_{\e\rightarrow 0} \be = 1-
      \frac{(f(2)-f(1))^2}{f(4)+f(1)-2f(2)} \\
  \gamma &= \lim_{\e\rightarrow 0} \gep =  \log_2\left(
  \frac{f(4)-f(2)}{f(2)-f(1)} \right)
\end{align*}
we must have $f(x) = ax^\gamma+b$.
\fi
\end{proof}

\end{document}
